% =====================================================================================
% Figure -- schematic plan of the study site, with overall extent.
% Belongs early in the methodology, where the site is first introduced, ahead of both
% the walkable-region figure (fig:walkable) and the marker figure (fig:markers).
%
% No extra packages required.
%
% Expected image file, relative to the main .tex:
%     figures/site-plan.png
%
% The source is portrait, roughly 2:3 (850 x 1290). At full \columnwidth in a two-column
% layout it stands about one and a half column widths tall, which is fine but dominant;
% 0.9\columnwidth is a safer default and is what is set below. If it must share a column
% with body text, drop to 0.8. Do NOT promote this to figure* -- stretched to full text
% width the plan becomes taller than the page.
%
% Export as PDF or SVG if the diagram tool allows it: the legend text and the parking
% stall hatching are fine detail that a low-resolution raster will lose.
%
% CAPTION CAVEAT -- two numbers to verify before submitting:
%   1. The 107 m x 160 m dimensions are read off the source diagram's own annotations.
%      Confirm they are the site extent and not, say, the scan bounding box.
%   2. The caption relates that extent to the 5854 m^2 walkable area cited in
%      figure_walkable.tex. 107 x 160 = 17,120 m^2 of bounding box, so the walkable
%      region is roughly a third of it. If the walkable figure's number changes, change
%      it here too, or drop the comparison sentence.
% Also confirm the site is de-identifiable enough for a double-blind submission -- a
% named campus plan can be recognisable even without a label.
%
% Reference in text as Figure~\ref{fig:site}.
% =====================================================================================

\begin{figure}[t]
  \centering
  \includegraphics[width=0.9\columnwidth]{figures/site-plan}
  \caption{Schematic plan of the study site, drawn to scale in an axonometric style with
    buildings extruded. The site spans roughly 107\,m by 160\,m. It is a mixed outdoor
    campus environment rather than an open field or a single corridor: a main building
    along the northern edge, a ring of smaller secondary structures, a service road, a
    surface parking area to the east, lawns of varying size, and a paved plaza with a
    circular planting bed at the southern end. This variety is deliberate. It gives the
    search task a range of sight lines, occluding geometry, and landmark types, and it
    supplies the durable concrete edges and corners on which the alignment markers are
    placed (Figure~\ref{fig:markers}). Participants do not traverse the whole extent
    shown: buildings, planting beds, the parking area, and the lawns are excluded, and
    the traversable region is the 5854\,m\textsuperscript{2} walkable surface of
    Figure~\ref{fig:walkable}.}
  \Description{A stylised top-down plan of a campus site in portrait orientation, drawn
    with buildings extruded slightly to show their height. Dimension arrows across the
    top and down the right side give the extent as about 107 metres wide by 160 metres
    tall. A legend in the lower left gives five categories by color: main building in
    pale pink, grass lawn in green, main road in grey, secondary structure in dark red,
    and circular plaza in pale blue-green. A large pale pink main building occupies the
    upper left. Dark red secondary structures are distributed around the edges of the
    site, including a row of long buildings along the right side and a row across the
    bottom. A grey road runs from the top of the plan down the middle and branches west.
    A long narrow green lawn runs down the centre of the site, with further green areas
    scattered throughout. A parking area with rows of marked stalls fills the upper
    right. At the bottom centre, a pale blue-green paved plaza contains a circular
    planting bed.}
  \label{fig:site}
\end{figure}
